% chapters/07/tujuan.tex
\begin{tujuanbox}
Bab ini melengkapi Blok~I dengan tabrakan yang benar, sehingga benda dapat
mendarat, terhalang, dan bersentuhan dengan benda lain. Setelah menuntaskan
bab ini, Anda mampu:
\begin{enumerate}
  \item menyusun benda permainan sebagai kelas turunan
        \code{pygame.sprite.Sprite} dan mengelolanya dengan
        \code{pygame.sprite.Group};
  \item memilih di antara \code{colliderect}, \code{spritecollide}, dan
        \code{groupcollide} sesuai keperluan;
  \item menyelesaikan tabrakan sumbu demi sumbu sehingga benda dapat mendarat
        di atas bongkahan sekaligus terhalang oleh sisinya;
  \item menjelaskan kapan persegi pembatas tidak memadai dan bagaimana
        \istilah{mask} dipakai sebagai gantinya;
  \item menganalisis tiga gejala tabrakan yang salah, yaitu lolos deteksi,
        tertangkap pada sisi yang salah, dan benda yang tertempel atau
        bergetar;
  \item memeriksa kebenaran penanganan tabrakan secara otomatis melalui
        program Anda sendiri.
\end{enumerate}
\end{tujuanbox}
